%\documentclass[12pt]{article}
\documentclass[12pt]{scrartcl}
\title{AY2026 MA05 Fourier Analysis Assignment 1}
\nonstopmode
%\usepackage[utf-8]{inputenc}
\usepackage{graphicx} % Required for including pictures
\usepackage[figurename=Figure]{caption}
\usepackage{float}    % For tables and other floats
\usepackage{verbatim} % For comments and other
\usepackage{amsmath}  % For math
\usepackage{amssymb}  % For more math
\usepackage{fullpage} % Set margins and place page numbers at bottom center
\usepackage{paralist} % paragraph spacing
\usepackage{listings} % For source code
\usepackage{subfig}   % For subfigures
%\usepackage{physics}  % for simplified dv, and
\usepackage{enumitem} % useful for itemization
\usepackage{siunitx}  % standardization of si units

\usepackage{tikz,bm} % Useful for drawing plots
%\usepackage{tikz-3dplot}
\usepackage{circuitikz}
\usepackage{hyperref} % For \url

%%% Colours used in field vectors and propagation direction
\definecolor{mycolor}{rgb}{1,0.2,0.3}
\definecolor{brightgreen}{rgb}{0.4, 1.0, 0.0}
\definecolor{britishracinggreen}{rgb}{0.0, 0.26, 0.15}
\definecolor{cadmiumgreen}{rgb}{0.0, 0.42, 0.24}
\definecolor{ceruleanblue}{rgb}{0.16, 0.32, 0.75}
\definecolor{darkelectricblue}{rgb}{0.33, 0.41, 0.47}
\definecolor{darkpowderblue}{rgb}{0.0, 0.2, 0.6}
\definecolor{darktangerine}{rgb}{1.0, 0.66, 0.07}
\definecolor{emerald}{rgb}{0.31, 0.78, 0.47}
\definecolor{palatinatepurple}{rgb}{0.41, 0.16, 0.38}
\definecolor{pastelviolet}{rgb}{0.8, 0.6, 0.79}
\begin{document}

\begin{center}
    \hrule
    \vspace{.4cm}
    {\textbf { \large AY2026 MA05 Fourier Analysis}}
\end{center}
% ====== TODO =======
%Please input your Name and Student ID below
{\textbf{Name:}\ Akito Shingo \hspace{\fill} \textbf{Deadline:} 17:00, Friday, May 22, 2026   \\
{ \textbf{Student ID:}} \ s1330096 \hspace{\fill} \textbf{Assignment:} 9 \\}
    \hrule

\begin{enumerate}

\item
Find the Discrete Fourier Transform by using Fast Fourier Transform for

$\{f(0),f(1),f(2),f(3)\}=\{0, 2, 3, 6\}$, $N=4$.

\textbf{Solution:} \\
% ====== TODO =======
% Please write your solution here.

The $N$-point DFT is defined by
\[
F(k)=\sum_{n=0}^{N-1} f(n)\,W_N^{\,nk},
\qquad W_N=e^{-2\pi i/N},\qquad k=0,1,\dots,N-1 .
\]
For $N=4$ the twiddle factor is $W_4=e^{-i\pi/2}=-i$, so
\[
W_4^{0}=1,\qquad W_4^{1}=-i,\qquad W_4^{2}=-1,\qquad W_4^{3}=i .
\]

\textbf{Step 1: split into even- and odd-indexed samples (decimation in time).}
\[
\text{even: }\{f(0),f(2)\}=\{0,3\},\qquad
\text{odd: }\{f(1),f(3)\}=\{2,6\}.
\]

\textbf{Step 2: the two 2-point DFTs.}
\begin{align*}
A(0)&=f(0)+f(2)=0+3=3, & A(1)&=f(0)-f(2)=0-3=-3,\\
B(0)&=f(1)+f(3)=2+6=8, & B(1)&=f(1)-f(3)=2-6=-4.
\end{align*}

\textbf{Step 3: combine with the butterfly} $F(k)=A(k)+W_4^{k}B(k)$,
$\;F(k+2)=A(k)-W_4^{k}B(k)$ for $k=0,1$:
\begin{align*}
F(0)&=A(0)+W_4^{0}B(0)=3+(1)(8)=11,\\
F(1)&=A(1)+W_4^{1}B(1)=-3+(-i)(-4)=-3+4i,\\
F(2)&=A(0)-W_4^{0}B(0)=3-(1)(8)=-5,\\
F(3)&=A(1)-W_4^{1}B(1)=-3-(-i)(-4)=-3-4i.
\end{align*}

\textbf{Result.}
\[
\boxed{\{F(0),F(1),F(2),F(3)\}=\{\,11,\ -3+4i,\ -5,\ -3-4i\,\}}
\]

\emph{Check (direct DFT).} $F(0)=0+2+3+6=11$, and
$F(1)=0+2(-i)+3(-1)+6(i)=-3+4i$, $F(2)=0+2(-1)+3(1)+6(-1)=-5$,
$F(3)=0+2(i)+3(-1)+6(-i)=-3-4i$, which agree with the FFT result.
Since $f$ is real, $F(3)=\overline{F(1)}$ as expected.

\\
Notice: \\
(1) Please submit your assignments according to the Homework Guidelines.

\url{https://web-ext.u-aizu.ac.jp/~xiangli/teaching/MA05/Guidelines.pdf}

(2) Write your Email title as "2026-A\{Assignment number\}-\{Your Student ID\}-\{Your Name\}".


\end{enumerate}

\end{document}
